\section{Kesimpulan}\label{sec:kesimpulan}

\subsection{Jawaban Rumusan Masalah}\label{sec:jawaban}

\textbf{RM1: mesin OCR mengubah cabang teks secara nyata.} Dengan embedding, arsitektur,
dan pembagian data dikunci, pilihan mesin OCR saja menggeser OA cabang teks 5,4 poin
(Tesseract-JPG 0,6970 sampai PaddleOCR 0,7514), konsisten di ketiga seed. Mesin modern
tidak otomatis unggul: PaddleOCR terbaik, EasyOCR di bawah teks acuan QS-OCR. Mesin juga
lebih menentukan daripada sumber citra, karena PaddleOCR pada JPG mengalahkan Tesseract
pada TIF asli. Proksi OOV tidak meramalkan urutan ini, sedangkan \textit{linear probe}
murah tanpa pelatihan meramalkannya.

\textbf{RM2: tidak dijawab.} BERT ditunda, sehingga interaksi mesin OCR dan embedding
belum diukur. Kesimpulan RM1 berlaku untuk embedding \textit{subword} statis.

\textbf{RM3: teks yang lebih baik terbawa sampai ke fusion, tetapi aturan rata-rata logit
tidak lebih tahan.} Dari Tesseract-JPG ke PaddleOCR, kedua aturan fusion naik sekitar
3 poin, dan sel terbaik (CONCAT dengan PaddleOCR, $0{,}8674 \pm 0{,}0014$) memanen 52\%
potensi oracle, di atas 43\% pada \textcite{audebert2020}. Untuk LATE, pemakaian teks naik
seiring akurasi cabang teks (Spearman 0,683); untuk CONCAT hubungannya lemah. LATE tidak
lebih tahan daripada CONCAT: rata-rata keduanya hampir sama, LATE unggul saat cabang citra
kuat, dan CONCAT unggul serta jauh lebih stabil saat cabang citra lemah.

\subsection{Koreksi atas Rancangan Sendiri}\label{sec:koreksi}

Uji \textit{missing modality} dengan pengganti kosong, yang kami rancang sebelum
eksekusi, ternyata hanya sah untuk LATE tanpa teks. Untuk CONCAT dan untuk semua versi
tanpa citra, pengganti kosong keluar distribusi dan membuat model jatuh di bawah model
unimodal modalitas yang tersisa. Uji lanjutan dengan pengganti rerata memastikan
penyebabnya, dan juga menggugurkan dugaan awal kami bahwa \textit{head} belajar
\enquote{teks kosong $\Rightarrow$ Note}. Pelajaran metodologisnya: pengganti dalam uji
\textit{missing modality} harus dipilih menurut aturan fusion yang diuji, dan kewajaran
hasilnya harus diperiksa terhadap model unimodal sebelum ditafsirkan.

\subsection{Keterbatasan}\label{sec:keterbatasan}

\begin{enumerate}
  \item \textbf{Tidak ada CER.} \enquote{Mesin lebih baik} di sini berarti teksnya lebih
  berguna untuk klasifikasi, bukan lebih benar dibaca.
  \item \textbf{Efek mesin bercampur dengan efek versi dan sumber citra.} Tesseract-JPG
  berbeda dari QS-OCR dalam versi (5.4 vs 4-beta) dan citra (JPG vs TIF).
  \item \textbf{Urutan baca EasyOCR dan PaddleOCR direkonstruksi heuristik} dengan asumsi
  satu kolom. Efeknya terhadap CNN~1D tidak diuji terpisah.
  \item \textbf{\textit{Head} CONCAT dilatih pada fitur latih} yang juga dipakai melatih
  model unimodal. Validasi silang di dalam data latih akan lebih bersih.
  \item \textbf{Tiga seed, dengan satu seed citra yang lemah} (0,7808). Sebagian
  kesimpulan tentang LATE vs CONCAT bertumpu pada seed itu.
  \item \textbf{Satu embedding dan satu dataset.} Hasil tidak dapat digeneralisasi ke BERT
  atau ke RVL-CDIP tanpa diuji.
\end{enumerate}

\subsection{Arah Lanjutan}\label{sec:arah}

\begin{enumerate}
  \item \textbf{Embedding kontekstual} untuk menjawab RM2, dengan pembagian data dan teks
  OCR yang sama.
  \item \textbf{Asal tarikan ke Note} pada CONCAT tanpa teks, misalnya dengan mengukur
  jarak fitur teks setiap kelas ke fitur teks kosong.
  \item \textbf{Validasi silang untuk \textit{head} CONCAT}, supaya \textit{head} menilai
  cabang dari perilaku di luar data latihnya.
\end{enumerate}

\subsection{Tahap Lanjutan: Teks sebagai Objek}\label{sec:lanjutan}

Sesuai arahan dosen pengampu, tahap berikutnya memperlakukan teks bukan hanya sebagai
rangkaian kata, tetapi sebagai \textbf{objek berposisi} pada halaman. Bahannya sudah
tersedia, karena setiap mesin OCR menyimpan kotak teksnya. Ada tiga kandidat
representasi:
\begin{enumerate}
  \item fitur posisi dua dimensi di samping embedding kata;
  \item peta kotak teks sebagai citra, yaitu setiap kotak digambar pada kanvas seukuran
  halaman;
  \item graf yang simpulnya kotak teks dan sisinya kedekatan spasial.
\end{enumerate}

Satu kontrol bersifat wajib. Kode ID arsip pada Tobacco3482 adalah objek teks berposisi
konsisten, dan fitur dari kode ID saja sudah memberi $\pm$60\% akurasi
\parencite{larson2025}. Model berbasis posisi dapat belajar jalan pintas ini alih-alih
belajar tata letak dokumen. Karena itu setiap model akan diuji dengan dan tanpa kotak
kode ID. Granularitas kotak juga berbeda antar-mesin: Tesseract per kata, sedangkan
EasyOCR dan PaddleOCR per baris. Rancangan rinci beserta referensinya akan diajukan
tersendiri.
